\chapter*{\center \Large \vspace{-4.5cm} ABSTRACT}
\addcontentsline{toc}{section}{ABSTRACT}
\markboth{ABSTRACT}{ABSTRACT} 
\thispagestyle{empty}
\begin{center}
\Large \vspace{-1.5cm} \textbf{NEXUSCARE: AI-BASED AUTOMATED AUDIT SOLUTION FOR MEDICAL CLAIMS}
\end{center}

\spacing{1.25}\noindent
The auditing of hospital medical claims in Peru is a predominantly manual process, consuming between 30 and 45 minutes per case. A typical batch of 101 claims requires 75.8 hours of work per specialized medical auditor, much of it devoted to routine, low-risk cases. This work presents NexusCare, a clinical audit agent that combines a deterministic contractual rules engine (Phase 1) with large language models (Phase 2) to classify each claim according to its clinical-economic risk and prioritize specialized review.

Model selection used a benchmark of five LLMs (Claude Sonnet 4, Claude Opus 4.6, GPT-5.4, Gemini 2.5 Pro, and MedGemma 27B) across 101 claims, totaling 505 evaluations. Sonnet 4 showed the highest discrimination capacity ($\sigma$ = 20.1), while other models exhibited over-penalization biases. A tiered strategy was adopted: Sonnet evaluates 100 \% of cases and escalates 33 \% to Opus, reducing cost by 57 \% compared to using Opus exclusively.

The integrated pipeline achieved 65.3 \% fast-track recommendations, with 66 low-risk cases validatable in approximately one minute, an API cost of USD 12.00, and 38.9 minutes of processing time. Threshold recalibration (+41.6 pp), anti-bias prompts (+6.0 pp), and contractual pre-filtering (-7.0 pp) account for the performance. The estimated reduction in total audit time ranges between 57 \% and 72 \%, exceeding the 40 \% target. Operational validation in production remains the next step.



\newpage
